\chapter{Entrelazamiento dodecafase--Witt y la lectura exacta de \texorpdfstring{\(3/11\)}{3/11}}
\label{chap:entrelazador-accion}
\index{entrelazador!dodecafase--Witt}
\index{Witt!módulo de incidencia}
\index{supervivencia!proyector de rango tres}

La aparición reiterada del cociente \(3/11\) en la pantalla dodecafásica,
en la polaridad de una estrella de Witt y en la filtración nonádica plantea
una pregunta más fuerte que la igualdad aritmética de tres fracciones. La
cuestión correcta es operatorial: determinar si esas tres apariciones son
trazas del mismo sector positivo visto en realizaciones diferentes y, si lo
son, construir los mapas que las entrelazan sin introducir un dato
metrológico posterior.

La respuesta posee cuatro niveles. Primero, la incidencia punto--hexada
construye una isometría canónica entre el módulo dodecafásico de dimensión
once y un autoespacio de dimensión once del diseño de Witt. Segundo, la
estrella de Witt actúa sobre ese mismo módulo, pero lo hace como involución
firmada; su \(3/11\) es un índice, no la traza de un proyector positivo.
Tercero, el sector orientado de supervivencia aparece en \(t=7\), donde un
único espacio de once ramas contiene un sector de rango tres. Cuarto, la
línea de orientación de ese triplete, junto con los testigos internos
\(B_K\) y \(u_K\), selecciona de manera variacional un único normalizador
\(S_3\), fija su orientación y produce por promedio de Reynolds y
descomposición polar el entrelazamiento isométrico sector--dodecafase. El
paso dodecafase--Witt es una segunda isometría de incidencia. En el calibre
etiquetado inicial ambas flechas no realizan todavía una única representación
de \(S_3\); la recalibración conjunta construida al final del capítulo transporta
la dodecada, el origen, la orientación, la cámara y la memoria, y cierra esa
equivarianza sin alterar el no-go del etiquetado fijo.
Los promedios de Reynolds, las descomposiciones en representaciones y el
factor polar utilizados en esta sección son los estándar
\cite{Serre1977Representations,HornJohnson2013}.

\begin{figure}[htbp]
\centering
\begin{tikzpicture}[
  font=\footnotesize,
  box/.style={draw,rounded corners=2mm,align=center,minimum width=2.35cm,minimum height=.98cm},
  arr/.style={-{Stealth[length=2.2mm]},thick},
  link/.style={<->,thick,dashed}
]
  \node[box,draw=hmtblue,fill=hmtlightblue] (seed) at (0,2.2)
    {semillas\\\(104\,976\)};
  \node[box,draw=hmtblue,fill=hmtlightblue] (u) at (3,2.2)
    {emisiones\\\(\mathcal U_6,\ |\mathcal U_6|=468\)};
  \node[box,draw=hmtgold,fill=hmtlightgold] (ppi) at (6,2.2)
    {microfibra\\\(\Pi_\pi^{(468)},\ \operatorname{rank}=5\)};
  \node[box,draw=hmtred,fill=hmtlightred] (d) at (9,2.2)
    {defecto positivo\\\(D_{\rm switch}\)};
  \node[box,draw=hmtviolet,fill=hmtlightviolet] (lam) at (12,2.2)
    {capacidad residual\\\(\lambda_*=1-\delta_*\)};

  \draw[arr,hmtblue] (seed)--node[above=16pt,font=\scriptsize,fill=white,inner sep=1pt]{\(F\)} (u);
  \draw[arr,hmtgold] (u)--node[above=16pt,font=\scriptsize,fill=white,inner sep=1pt]{\(\tau_{468}\)} (ppi);
  \draw[arr,hmtred] (ppi)--node[above=16pt,font=\scriptsize,fill=white,inner sep=1pt]{\(D_{\rm switch}\)} (d);
  \draw[arr,hmtviolet] (d)--node[above=16pt,font=\scriptsize,fill=white,inner sep=1pt]{\(\tau+\mathrm{CR}\)} (lam);

  \node[box,draw=hmtgold,fill=hmtlightgold] (g) at (1.5,-.2)
    {sector orientado \(t=7\)\\\(L_{\rm or}\otimes(H_G,P_G)\)};
  \node[box,draw=hmtgreen,fill=hmtlightgreen] (v) at (5,-.2)
    {dodecafase\\\(V_{11},P_3\)};
  \node[box,draw=hmtgreen,fill=hmtlightgreen] (e) at (8.5,-.2)
    {incidencia Witt\\\(E_{30},Q_3\)};
  \node[box,draw=hmtgray,fill=hmtpaper] (t) at (12,-.2)
    {lector espectral\\\(\mathcal L_{\beta,*}\to T_{\rm band}\)};

  \draw[arr,hmtgold] (g)--node[above=16pt,font=\scriptsize,fill=white,inner sep=1pt]{\(W_{GK}\)} (v);
  \draw[arr,hmtgreen] (v)--node[above=16pt,font=\scriptsize,fill=white,inner sep=1pt]{\(U_W=I/6\)} (e);
  \draw[arr,hmtgray] (lam)--(t);
  \draw[arr,hmtred] (v) to[bend left=20]
    node[below=3pt,font=\scriptsize,fill=white,inner sep=1pt]{\((3/11)\Delta_4\)} (d);
\end{tikzpicture}

\caption{Arquitectura del cierre pre-dimensional. La fila superior construye
la medida y la acción desde el cociente de emisiones; la inferior transporta
el sector orientado \(t=7\) hacia la dodecafase y, después, la dodecafase hacia
Witt. Las flechas continuas son isometrías de proyectores demostradas. La
primera se fija por la selección variacional del
\cref{thm:identificacion-isometrica-tres-proyectores}; la segunda no es
\(H_*\)-equivariante en el calibre fijo y se vuelve equivariante mediante el
cambio de calibre íntegro del \cref{thm:recalibracion-hstar-witt}.}
\label{fig:entrelazador-accion}
\end{figure}

\section{Tipado de los objetos asociados a \texorpdfstring{\(3/11\)}{3/11}}

Sea
\[
 V_{12}=\mathbb R^{12},
 \qquad
 V_{11}=\mathbf 1^\perp
 =\left\{v\in\mathbb R^{12}:\sum_{p=1}^{12}v_p=0\right\}.
\]
La acción icosaédrica declarada sobre las doce posiciones se descompone como
\[
 V_{12}=V_1\oplus V_3\oplus V_{3'}\oplus V_5.
\]
Denotemos por \(P_3\) el proyector ortogonal sobre el primer sumando
tridimensional. Entonces
\[
 P_3^2=P_3=P_3^*,
 \qquad
 \operatorname{rank}P_3=3,
 \qquad
 \tau_{11}(P_3)
 =\frac1{11}\operatorname{Tr}_{V_{11}}P_3
 =\frac3{11}.
\]

Hay tres objetos próximos y un cuarto objeto de tipo distinto:
\begin{enumerate}
  \item la incidencia punto--hexada transporta \(P_3\) a un proyector
  \(Q_3\) de rango tres sobre un módulo de Witt;
  \item la supervivencia en \(t=7\) distingue tres ramas entre once y
  produce un proyector \(P_G\) con traza normalizada \(3/11\);
  \item un entrelazador orientado identifica
  \(P_3V_{11}\) con la realización orientada de \(P_GH_G\);
  \item una estrella de Witt tiene siete modos pares y cuatro impares, por
  lo que su traza normalizada también vale \(3/11\), pero contiene espectro
  negativo y no es un proyector.
\end{enumerate}
La distinción entre proyector e involución será esencial para que la ley de
acción de la sección siguiente sea positiva.

\section{La matriz de incidencia punto--hexada}

Sea \(\mathcal H\) el conjunto de las \(132\) hexadas del diseño
\[
 W_{12}=S(5,6,12).
\]
Definimos
\[
 I:\mathbb R^{12}\longrightarrow\mathbb R^{\mathcal H},
 \qquad
 (Iv)(H)=\sum_{p\in H}v_p,
\]
o, equivalentemente,
\[
 I_{H,p}=\mathbf 1_{\{p\in H\}}.
\]
Este operador sólo registra incidencia. No contiene \(\alpha\), unidades,
masas, CKM, Barbero ni ningún objetivo físico.

\begin{lemma}[Identidad métrica de la incidencia]
\label{lem:incidencia-metrica-witt}
Se cumple
\[
 \boxed{I^*I=36I_{12}+30J_{12}.}
\]
Por tanto, para todo \(v\in V_{11}\),
\[
 \|Iv\|^2=36\|v\|^2.
\]
\end{lemma}

\begin{proof}
En un sistema \(S(5,6,12)\), cada punto pertenece a
\[
 r=\frac{\binom{11}{4}}{\binom54}=66
\]
hexadas, y cada par de puntos pertenece a
\[
 \lambda_2=\frac{\binom{10}{3}}{\binom43}=30
\]
hexadas. Las entradas diagonales de \(I^*I\) valen \(66\) y las no
diagonales valen \(30\). Luego
\[
 I^*I=(66-30)I_{12}+30J_{12}.
\]
Como \(J_{12}\) se anula sobre \(\mathbf1^\perp\), se obtiene la segunda
identidad.
\end{proof}

Para identificar la rama espectral seleccionada por \(I\), introducimos
\[
 G_{H,H'}=\binom{|H\cap H'|}{4}.
\]

\begin{lemma}[Identidad espectral de la incidencia]
\label{lem:incidencia-espectral-witt}
Con \(J_{132,12}\) la matriz rectangular de unos,
\[
 \boxed{GI=30I+15J_{132,12}.}
\]
En particular, si \(v\in V_{11}\), entonces
\[
 G(Iv)=30Iv.
\]
\end{lemma}

\begin{proof}
La entrada \((H,p)\) de \(GI\) es
\[
 \sum_{H'\ni p}\binom{|H\cap H'|}{4}.
\]
El censo exhaustivo de \(W_{12}\) da \(45\) cuando \(p\in H\) y \(15\)
cuando \(p\notin H\). Por consiguiente, la entrada vale
\(30I_{H,p}+15\). La reconstrucción de las \(132\) hexadas desde el código
de Golay ternario extendido comprueba las \(132\cdot12\) entradas como
enteros.
\end{proof}

Definamos
\[
 E_{30}^{\rm inc}:=I(V_{11})\subseteq\ker(G-30I).
\]
Su dimensión es once por el \cref{lem:incidencia-metrica-witt}. El cálculo
módulo \(1009\) obtiene
\[
 \operatorname{rank}_{\mathbb F_{1009}}(G-30I)=121.
\]
Como la reducción módulo \(1009\) da una cota inferior para el rango
racional, se sigue
\(\dim_{\mathbb Q}\ker(G-30I)\le11\). La inclusión anterior proporciona
once vectores linealmente independientes en ese núcleo. Por tanto,
\[
 E_{30}^{\rm inc}=\ker(G-30I).
\]
En adelante escribiremos \(E_{30}\).

\begin{theorem}[Entrelazador canónico dodecafase--Witt]
\label{thm:entrelazador-dodecafase-witt}
El operador
\[
 \boxed{
 U_W=\frac16 I\big|_{V_{11}}:V_{11}\longrightarrow E_{30}
 }
\]
es una isometría sobreyectiva. Para todo
\(g\in\operatorname{Aut}(W_{12})\cong M_{12}\),
\[
 U_W\rho_{12}(g)=\rho_{\mathcal H}(g)U_W.
\]
Entre los kernels punto--hexada equivariantes, la condición de isometría y
el signo positivo sobre las incidencias determinan \(U_W\) de manera única.
\end{theorem}

\begin{proof}
La isometría se sigue de \(I^*I=36I\) sobre \(V_{11}\). Su imagen tiene
dimensión once y está contenida en \(E_{30}\), también de dimensión once;
por tanto es sobreyectiva. La equivariancia traduce
\[
 p\in H\quad\Longleftrightarrow\quad g(p)\in g(H).
\]

Para la unicidad, \(M_{12}\) es transitivo sobre las banderas
\((p,H)\), con \(p\in H\), y sobre las antibanderas. Todo núcleo
equivariante tiene dos valores, \(a\) sobre banderas y \(b\) sobre
antibanderas:
\[
 T=aI+b(J-I)=(a-b)I+bJ.
\]
El término \(J\) se anula sobre \(V_{11}\), de modo que allí \(T\) es un
múltiplo de \(I\). La isometría fuerza \(|a-b|=1/6\) y la normalización
positiva selecciona \(a-b=1/6\).
\end{proof}

\section{Transporte positivo del proyector trítico}

El entrelazador permite definir
\[
 \boxed{Q_3=U_WP_3U_W^*.}
\]
Se obtiene
\[
 Q_3^2=Q_3=Q_3^*,
 \qquad
 \operatorname{rank}Q_3=3,
\]
y el diagrama
\[
\begin{CD}
V_{11} @>{U_W}>> E_{30}\\
@V{P_3}VV @VV{Q_3}V\\
V_{11} @>{U_W}>> E_{30}
\end{CD}
\]
conmuta. En particular,
\[
 \boxed{\tau_{V_{11}}(P_3)=\tau_{E_{30}}(Q_3)=\frac3{11}.}
\]

La canonicidad de \(U_W\) pertenece al diseño etiquetado y a su acción de
automorfismos. Debe distinguirse de la pantalla \(A_5/C_5\) usada para
definir \(P_3\): en la realización etiquetada construida, sólo la identidad
de esa acción concreta de \(A_5\) preserva el conjunto fijo de hexadas. Por
ello \(U_W\) transporta canónicamente el \(P_3\) ya dado hacia
\(Q_3\), pero no convierte por sí solo la acción \(A_5/C_5\) en una
subacción de \(M_{12}\). El puente con el sector orientado se construirá por separado
desde los datos marcados \(K\) y \(B_K\).

\section{Índice firmado de la estrella de Witt}

Sea \(B\) una tetrada. Las cuatro hexadas que la contienen se escriben
\[
 H_i=B\sqcup P_i,\qquad |P_i|=2,
\]
donde los cuatro pétalos \(P_i\) particionan el complemento de \(B\).
La involución de estrella \(R_B\) fija \(B\) e intercambia los dos puntos
de cada pétalo. Su tipo de ciclos sobre doce puntos es \(1^4\,2^4\). En
\(V_{11}\), y mediante \(U_W\) también en \(E_{30}\), su espectro es
\[
 \operatorname{Spec}(R_B)=\{1^{[7]},(-1)^{[4]}\}.
\]
Por ello,
\[
 \boxed{\tau_{11}(R_B)=\frac{7-4}{11}=\frac3{11}.}
\]

\begin{theorem}[Obstrucción proyector--estrella]
\label{thm:no-go-proyector-estrella}
No existe un isomorfismo \(F:V_{11}\to E_{30}\) tal que
\[
 FP_3=R_BF.
\]
\end{theorem}

\begin{proof}
Si existiera, \(P_3\) y \(R_B\) serían semejantes. Sin embargo,
\[
 \operatorname{Spec}(P_3)=\{1^{[3]},0^{[8]}\},
 \qquad
 \operatorname{Spec}(R_B)=\{1^{[7]},(-1)^{[4]}\}.
\]
Sus polinomios mínimos son \(x(x-1)\) y
\((x-1)(x+1)\), respectivamente. La semejanza es imposible.
\end{proof}

La igualdad \(3/11\) posee así dos lecturas sobre el mismo módulo. Para
\(Q_3\) es la masa de un sector positivo de rango tres. Para \(R_B\) es
un índice de paridad. En una acción positiva debe intervenir \(Q_3\); la
estrella conserva orientación y distingue hojas, pero no reemplaza al
proyector.

\section{Reducción homogénea \texorpdfstring{\(11\to3\to1\)}{11-3-1} en \texorpdfstring{\(t=7\)}{t=7}}

La lectura anterior de \(3/11\) en \(t=8\) comparaba tres raíces locales
con once continuaciones de una sola raíz. Numerador y denominador no
pertenecían al mismo espacio. La tabla detallada localiza el objeto correcto
en \(t=7\): hay once ramas locales; las once sobreviven un horizonte, tres
sobreviven dos y una sobrevive tres,
\[
 \boxed{11\longrightarrow3\longrightarrow1.}
\]

\begin{center}
\small
\begin{tabularx}{\textwidth}{cYcc}
\toprule
índice&triple de palabras&\(h=2\)&\(h=3\)\\
\midrule
0&\texttt{200112|212221|110110}&3&0\\
1&\texttt{200121|212212|110110}&3&1\\
2&\texttt{200211|212122|110110}&3&0\\
3&\texttt{202122|210211|110110}&0&0\\
4&\texttt{202212|210121|110110}&0&0\\
5&\texttt{202221|210112|110110}&0&0\\
6&\texttt{210111|202222|110110}&0&0\\
7&\texttt{210222|202111|110110}&0&0\\
8&\texttt{212112|200221|110110}&0&0\\
9&\texttt{212121|200212|110110}&0&0\\
10&\texttt{212211|200122|110110}&0&0\\
\bottomrule
\end{tabularx}
\end{center}

Sea \(\mathcal B_7\) el conjunto de estas once ramas y
\[
 H_G=\ell^2(\mathcal B_7).
\]
Definimos
\[
 P_G\delta_b=
 \begin{cases}
 \delta_b,&b\text{ posee continuación a horizonte dos},\\
 0,&\text{en otro caso}.
 \end{cases}
\]
Entonces
\[
 P_G^2=P_G=P_G^*,
 \qquad
 \operatorname{rank}P_G=3,
 \qquad
 \boxed{\tau_{H_G}(P_G)=\frac3{11}.}
\]
Ésta es la tercera realización positiva.

\section{La acción residual de \texorpdfstring{\(S_3\)}{S3} y las regiones positivas de \texorpdfstring{\(\pi\)}{pi}}

Una permutación de las tres posiciones finales, aplicada simultáneamente a
las dos primeras palabras de cada rama y dejando fija la tercera, preserva
\(\mathcal B_7\). Esta regla define una acción de \(S_3\). Sus órbitas
tienen tamaños
\[
 3+3+1+1+3=11.
\]
El soporte de \(P_G\), formado por los índices \(0,1,2\), es una órbita
completa, por lo que \(P_G\) conmuta con \(S_3\).

El carácter de la representación sobre las once ramas, leído en identidad,
transposición y ciclo de orden tres, es
\[
 \chi_{\mathcal B_7}=(11,5,2)
 =5\chi_{\rm triv}+3\chi_{\rm std}.
\]
Sobre la imagen de \(P_G\),
\[
 \chi_G=(3,1,0)
 =\chi_{\rm triv}+\chi_{\rm std}.
\]

Las tres regiones positivas de la microfibra de \(\pi\) tienen firmas
\[
 339555,\qquad393555,\qquad933555.
\]
Las colas de la primera palabra de las tres ramas supervivientes son
\[
 112,\qquad121,\qquad211.
\]
La sustitución coordenada \(1\mapsto3\), \(2\mapsto9\) da
\[
\begin{array}{c|c|c}
\text{rama}&\text{cola}&\text{región positiva}\\ \hline
0&112&339555\\
1&121&393555\\
2&211&933555.
\end{array}
\]

\begin{theorem}[Puente horizontal--vertical de la microfibra positiva]
\label{thm:puente-pi-puerta-s3}
La regla anterior es una biyección \(S_3\)-equivariante entre las tres
regiones positivas de \(\pi\) y las tres ramas seleccionadas por \(P_G\).
El mapa conserva la posición del único símbolo excepcional.
\end{theorem}

\begin{proof}
Ambos conjuntos son las tres palabras obtenidas al colocar un único símbolo
excepcional en una de tres posiciones. La sustitución se aplica coordenada a
coordenada y, por tanto, conmuta con toda permutación de posiciones. La
posición excepcional recupera unívocamente el elemento de origen.
\end{proof}

\section{Dependencia de la orientación en el entrelazador dodecafásico--Witt}

Para cualquier normalizador \(H=N_{A_5}(C_3)\simeq S_3\) de un subgrupo de
orden tres en \(A_5\), definimos el carácter abstracto del cociente
\[
 \operatorname{sgn}_H:H\longrightarrow H/C_3\simeq C_2
 \longrightarrow\{+1,-1\},
 \qquad
 \operatorname{sgn}_H(h)=
 \begin{cases}
 +1,&h\in C_3,\\
 -1,&h\in H\setminus C_3.
 \end{cases}
\]
Este carácter no es el signo ambiental de una permutación de cinco letras:
todo elemento de \(A_5\) tiene signo ambiental \(+1\). Es el carácter
intrínseco de la reflexión en el cociente \(H/C_3\). Con esta convención, el
carácter del sector icosaédrico restringido a \(H\) es
\[
 \chi_{P_3V_{11}|H}=(3,-1,0)
 =\chi_{\operatorname{sgn}_H}+\chi_{\rm std}.
\]
En cambio,
\[
 \chi_{\operatorname{im}P_G}=(3,1,0)
 =\chi_{\rm triv}+\chi_{\rm std}.
\]

\begin{corollary}[Obstrucción al entrelazador no orientado]
\label{cor:obstruccion-entrelazador-incidencia-p3}
No existe un isomorfismo \(S_3\)-equivariante
\[
 \operatorname{im}P_G\longrightarrow P_3V_{11}
\]
para las acciones no retorcidas.
\end{corollary}

\begin{proof}
Representaciones isomorfas poseen el mismo carácter. Los valores en una
transposición son \(1\) y \(-1\).
\end{proof}

Sea \((e_0,e_1,e_2)\) la base posicional de
\(\mathbb R^3_{\rm perm}\). Definimos
\[
 L_{\rm or}
 =\bigwedge\nolimits^3\mathbb R^3_{\rm perm},
 \qquad
 \ell_{\rm or}=e_0\wedge e_1\wedge e_2
\]
la línea de orientación del triplete, provista del ancla positiva
\(\ell_{\rm or}\). La acción de \(S_3\) sobre esta línea es precisamente su
carácter de signo abstracto. Al tensorizar por ella,
\[
 \chi_{L_{\rm or}\otimes\operatorname{im}P_G}
 =(3,-1,0),
\]
de modo que la obstrucción de caracteres desaparece. La igualdad de
caracteres demuestra la existencia de una clase de isomorfismos, pero aún no
selecciona uno. La selección se obtiene de dos testigos ya presentes en la
dodecafase.

\section{Selección variacional del normalizador \texorpdfstring{\(S_3\)}{S3}}

Sea
\[
 K=(234,543,140,729,659,824,621,58,914,794,146,601)
\]
la carta decimal del sello. El testigo común mínimo de las dos lecturas es
\[
 B_K=H_e\cap P_\pi
 =\{p:K_p\ge729\}
 =\{4,6,9,10\},
\]
en la numeración matemática \(1,\ldots,12\). En la convención equivalente
con índices desde cero, el mismo conjunto es \(\{3,5,8,9\}\). Definimos dos
vectores del sector \(P_3\):
\[
 b_K=P_3\left(\mathbf1_{B_K}-\frac{|B_K|}{12}\mathbf1\right),
 \qquad
 u_K=P_3\left(K-\overline K\,\mathbf1\right).
\]
El primero satisface exactamente
\[
 \|b_K\|^2=1.
\]

El grupo \(A_5\) contiene diez subgrupos \(C_3\). Para cada uno, sea
\[
 H(C_3)=N_{A_5}(C_3)\simeq S_3
\]
y sea el proyector de signo
\[
 e_{{\rm sgn},H}
 =\frac16\sum_{h\in H}\operatorname{sgn}_H(h)\rho(h).
\]
Aquí \(\operatorname{sgn}_H\) es el carácter abstracto de \(H/C_3\)
definido arriba, no el signo ambiental —trivial— de \(A_5\).
La energía de orientación asociada al testigo común es
\[
 \mathcal E_B(H)=\|e_{{\rm sgn},H}b_K\|^2.
\]
El censo exacto en \(\mathbb Q(\sqrt5)\) produce un máximo único:
\[
 \mathcal E_B(H_*)
 =\frac23+\frac{2}{15}\sqrt5.
\]
El segundo nivel es
\[
 \frac13+\frac{2}{15}\sqrt5,
\]
de modo que la brecha exacta es
\[
 \boxed{\mathcal E_B(H_*)-\mathcal E_B(H_{\rm segundo})=\frac13.}
\]
El mismo \(H_*\) es el máximo único al reemplazar \(b_K\) por \(u_K\).
Por tanto, el selector no depende de una sola codificación del testigo.

En esta etiquetación,
\[
 H_*=N_{A_5}\bigl(\langle c_*\rangle\bigr),
\qquad
 c_*=(2\,3\,4).
\]
La regla es natural bajo relabelización simultánea por \(A_5\): al
transportar \(B_K\), \(K\) y \(P_3\), el máximo se conjuga por el mismo
elemento.

\section{Selección de orientación y presentación de \texorpdfstring{\(S_3\)}{S3}}

Dentro de \(H_*\), consideremos
\[
 s(g)=\langle b_K,\rho(g)u_K\rangle.
\]
El máximo entre los dos elementos de orden tres selecciona unívocamente
\[
 c_*=(0,1,3,4,2),
\]
en notación de imágenes de la permutación de cinco puntos, y el máximo entre
las tres involuciones selecciona
\[
 r_*=(1,0,4,3,2).
\]
Se verifica exactamente
\[
 r_*c_*r_*=c_*^{-1}.
\]
La presentación del sector orientado mediante el ciclo de posiciones
\(c=(0\,1\,2)\) y la reflexión \(r=(0\,2)\), que fija la posición central,
determina entonces un único isomorfismo
\[
 \theta:S_3^{\rm or}\xrightarrow{\sim}H_*,
 \qquad
 \theta(c)=c_*,
 \qquad
 \theta(r)=r_*.
\]
Aquí \(S_3^{\rm or}\) denota el grupo de reetiquetados del sector orientado
\(t=7\), generado por \(c\) y \(r\); no introduce una clase operatoria
adicional.
La única rama que alcanza horizonte tres es precisamente la rama central
\(1\); este dato fija cuál de las tres posiciones queda inmóvil bajo la
reflexión.

\section{Promedio de Reynolds y factor polar}

Sea
\[
 W_{\rm or}=L_{\rm or}\otimes\mathbb R^3_{\rm perm}
\]
el triplete del sector orientado. Con el ancla ya fijada, tomamos
\[
 \boxed{w=\ell_{\rm or}\otimes e_1,}
\]
donde \(e_1\) representa la posición central, la única que alcanza horizonte
tres. Definimos
\[
 A
 =\sum_{g\in S_3}
 \rho_{W_{\rm or}}(g)\,
 |w\rangle\langle b_K|\,
 \rho_{V_{11}}(\theta(g))^{-1}.
\]
Por construcción,
\[
 A\rho_{V_{11}}(\theta(g))
 =\rho_{W_{\rm or}}(g)A,
\qquad
 A=AP_3.
\]
El cálculo exacto en \(\mathbb Q(\sqrt5)\) obtiene
\[
 \operatorname{rank}A=3.
\]
Además, \(G_A=AA^*\) tiene todas sus entradas diagonales iguales a
\[
 3+\frac25\sqrt5
\]
y todas sus entradas no diagonales iguales a
\[
 \frac52+\frac35\sqrt5.
\]
Sus dos autovalores isótipo son
\[
 \lambda_{\rm sgn}=8+\frac85\sqrt5,
 \qquad
 \lambda_{\rm std}=\frac12-\frac15\sqrt5.
\]
Ambos son estrictamente positivos. En particular, \(G_A^{-1/2}\) existe y
la descomposición polar
\[
 \boxed{U_{GK}=G_A^{-1/2}A}
\]
es una isometría equivariante de \(P_3V_{11}\) sobre \(W_{\rm or}\):
\[
 U_{GK}^*U_{GK}=P_3,
 \qquad
 U_{GK}U_{GK}^*=I_{W_{\rm or}}.
\]

Sea
\[
 J_G:W_{\rm or}\longrightarrow L_{\rm or}\otimes H_G
\]
la inclusión isométrica en las tres ramas seleccionadas. Entonces
\[
 J_GJ_G^*=I_{L_{\rm or}}\otimes P_G.
\]
La isometría parcial sector--dodecafase se denota
\[
 \boxed{
 W_{GK}=U_{GK}^*J_G^*
 :L_{\rm or}\otimes H_G\longrightarrow V_{11}
 }
\]
satisface
\[
 W_{GK}^*W_{GK}
 =I_{L_{\rm or}}\otimes P_G,
\qquad
 W_{GK}W_{GK}^*=P_3,
\]
y, por tanto,
\[
 P_3W_{GK}
 =W_{GK}
 =W_{GK}(I_{L_{\rm or}}\otimes P_G).
\]

\begin{theorem}[Identificación isométrica relativa de tres proyectores]
\label{thm:identificacion-isometrica-tres-proyectores}
Relativamente a la acción \(A_5/C_5\), al proyector \(P_3\), al sello
\(K\), al testigo doblemente construido \(B_K\), al sector finito orientado
\(t=7\), al carácter abstracto \(\operatorname{sgn}_{H_*}\), al ancla
\(\ell_{\rm or}\) y a la regla variacional \(\mathcal E_B\), existe una
única isometría parcial normalizada \(W_{GK}\) que identifica \(P_G\) con
\(P_3\). Independientemente, la isometría de incidencia \(U_W\) identifica
\(P_3\) con \(Q_3\). En consecuencia, se obtiene la siguiente cadena de
identificaciones isométricas de proyectores:
\[
\begin{CD}
L_{\rm or}\otimes H_G
 @>{W_{GK}}>>
 V_{11}
 @>{U_W}>>
 E_{30}\\
@V{I_{L_{\rm or}}\otimes P_G}VV
 @VV{P_3}V
 @VV{Q_3}V\\
L_{\rm or}\otimes H_G
 @>{W_{GK}}>>
 V_{11}
 @>{U_W}>>
 E_{30}.
\end{CD}
\]
Las trazas tipadas satisfacen
\[
 \tau_{L_{\rm or}\otimes H_G}(I_{L_{\rm or}}\otimes P_G)
 =\tau_{V_{11}}(P_3)
 =\tau_{E_{30}}(Q_3)
 =\frac3{11}.
\]
El enunciado, en el calibre fijo, identifica isométricamente tres
proyectores; no afirma todavía que sus tres espacios sean una misma
representación de un grupo común.
\end{theorem}

\begin{proof}
La energía \(\mathcal E_B\) selecciona un único normalizador \(H_*\);
el criterio de puntuación \(s(g)\) selecciona un único par \((c_*,r_*)\), y con él un
único isomorfismo \(\theta\). El promedio de Reynolds es entonces
equivariante para las acciones de \(S_3^{\rm or}\) y \(H_*\) relacionadas
por \(\theta\), y posee rango tres. La parte positiva de su descomposición
polar fija la normalización sin introducir una base. Las identidades
parciales anteriores demuestran que \(W_{GK}\) identifica
\(P_G\) y \(P_3\). Por otra parte, el
\cref{thm:entrelazador-dodecafase-witt} identifica \(P_3\) y \(Q_3\)
como proyectores mediante \(U_W\). Estas dos identidades operatoriales hacen
conmutar ese diagrama de calibre fijo, sin extender allí la equivarianza de la
primera flecha a la segunda.
\end{proof}

\begin{proposition}[Obstrucción a la equivarianza triple en el calibre fijado]
\label{prop:no-go-equivarianza-hstar-witt}
La construcción anterior no demuestra que \(U_W\) sea
\(H_*\)-equivariante. En el etiquetado declarado, la acción fija de
\(A_5/C_5\) sobre los doce puntos preserva el conjunto de hexadas de
\(W_{12}\) únicamente para el elemento identidad. Por tanto, el
entrelazamiento \(H_*\)-equivariante de \(W_{GK}\) no se prolonga, con ese
calibre, a una representación común sobre \(E_{30}\).
\end{proposition}

\begin{proof}
La primera afirmación es una separación de dominios de equivarianza:
\(W_{GK}\) usa \(H_*\subset A_5\), mientras que el teorema de incidencia
prueba la equivarianza de \(U_W\) para
\(\operatorname{Aut}(W_{12})\). El censo exacto reproducible aplica los sesenta
elementos de la acción \(A_5/C_5\) al sistema de hexadas etiquetado; sólo la
identidad conserva ese sistema. Como \(H_*\) contiene seis elementos, su
acción no puede entrar por restricción de esa acción fija en
\(\operatorname{Aut}(W_{12})\).
\end{proof}

La proposición anterior aísla exactamente el calibre fijo; no es una
obstrucción abstracta a transportar conjuntamente la calibración. El
propietario rector de recalibración \(S_3\)--Witt construye el monomorfismo y
la segunda flecha equivariantes como sigue.

\begin{theorem}[Recalibración conjunta \texorpdfstring{\(H_*\)}{H*}--Witt]
\label{thm:recalibracion-hstar-witt}
Sea
\(\bar\rho_{\rm src}:H_*\to S_{12}\) la acción puntual de las doce clases
laterales determinada por \(c_*\), \(r_*\) y el isomorfismo \(\theta\) del
sector orientado, y sea \(\rho_{\rm src}\) su linearización en \(V_{11}\).
En convención de imágenes uno-basada, la permutación
\[
 \boxed{\sigma=(1,5,2,6,7,4,3,10,11,9,8,12)}
\]
define el monomorfismo
\[
 \boxed{
 \iota_\sigma:H_*\hookrightarrow\operatorname{Aut}(W_{12})\cong M_{12},
 \qquad
 \iota_\sigma(h)=
 \sigma\,\bar\rho_{\rm src}(h)\,\sigma^{-1}.
 }
\]
Si se transporta íntegramente la calibración mediante
\[
 W_{12}^{\sigma}=\{\sigma^{-1}B:B\in W_{12}\},
 \qquad
 \widetilde U_W:=U_WP_\sigma,
\]
donde \(P_\sigma\) es la unitaria de permutación, entonces, bajo la
identificación correspondiente del espacio de incidencia transportado con
\(E_{30}\), para todo \(h\in H_*\),
\[
 \boxed{
 \widetilde U_W\,\rho_{\rm src}(h)
 =\rho_{E_{30}}\!\bigl(\iota_\sigma(h)\bigr)\,\widetilde U_W.
 }
\]
En particular,
\[
 \widetilde Q_3=\widetilde U_WP_3\widetilde U_W^*
\]
es la realización Witt del mismo proyector y la cadena
\[
 I_{L_{\rm or}}\otimes P_G
 \xleftrightarrow{\ W_{GK}\ }
 P_3
 \xleftrightarrow{\ \widetilde U_W\ }
 \widetilde Q_3
\]
porta, tras identificar la acción fuente mediante \(\theta\), una
representación común de \(H_*\), con la acción de llegada dada por
\(\iota_\sigma\). La recalibración transporta conjuntamente
\(C_W,W_{12},I\), la dodecada, el origen, la orientación, la cámara y la
memoria; no consiste en sustituir aisladamente \(U_W\).
\end{theorem}

\begin{proof}
El certificado exacto reconstruye el código de Golay ternario, las \(132\)
hexadas, el grupo \(M_{12}\) de orden \(95040\), el \(S_3\) fuente y un
\(S_3\) semirregular de llegada, ambos de orden seis. Comprueba para cada
\(h\in H_*\) que
\(
\sigma\bar\rho_{\rm src}(h)\sigma^{-1}\in M_{12}
\), que \(W_{12}^{\sigma}\) es \(H_*\)-invariante y que su incidencia
punto--hexada es equivariante. La inyectividad de \(\iota_\sigma\) se sigue
de que es conjugación de la acción fiel fuente. Además,
\[
 P_\sigma\rho_{\rm src}(h)
 =\rho_{12}\!\bigl(\iota_\sigma(h)\bigr)P_\sigma.
\]
Al componer esta identidad con la equivarianza \(M_{12}\) de \(U_W\) se
obtiene la identidad enmarcada. Como \(P_3\) conmuta con \(H_*\), su
transporte \(\widetilde Q_3\) conmuta con \(\iota_\sigma(H_*)\); la
equivarianza de \(W_{GK}\) ya demostrada completa la composición. El mismo
certificado verifica que el diseño original no es invariante bajo la acción
fuente, de modo que se conserva, sin contradicción, el no-go del calibre fijo.
\end{proof}

\textsc{Propietario material.} Gobiernan este resultado el teorema
\textsc{HMT--RDF--EXC--S3}, secciones 3.1--3.3 del delta rector de propagación
excepcional Witt--Moonshine, y su certificado reproducible
\texttt{verificar\_recalibracion\_s3\_witt.py}, cuya salida es
\texttt{PASS\_RECALIBRACION\_S3\_WITT}. El estatuto es
\texttt{FORMALIZACION\_NUEVA/CERTIFICADO\_NUEVO}; no se postula unicidad del
cambio de calibre.

Antes de usar \(B_K\), \(u_K\), la orientación y el factor polar había
\(10\cdot6\cdot4=240\) mapas etiquetados, o \(40\) después de identificar
relabelizaciones internas del sector. Las reglas anteriores reducen ese
conjunto a uno. Éste es el contenido verificable de la palabra
\emph{canónico} para \(W_{GK}\), dentro de las reglas declaradas. La acción
común sobre el módulo de Witt procede del cambio explícito \(\sigma\); no se
afirma que ese cambio de calibre sea único o canónico.

La regla variacional \(\mathcal E_B\) utiliza únicamente datos HMT anteriores
y fija la clase en la que se obtiene la unicidad. No identifica la involución
de estrella con \(P_3\), pues sus espectros siguen siendo distintos.

\section{Comparación de las lecturas de \texorpdfstring{\(t=7\) y \(t=8\)}{t=7 y t=8}}

En \(t=8\) hay tres raíces locales y once extensiones a horizonte nueve de
la única raíz que continúa. Los tamaños de las tres fibras son
\[
 (11,0,0).
\]
La matriz natural rama--cadena tiene una fila de once unos y dos filas
nulas, por lo que su rango es uno. En \(t=9\), además, el cociente visible
cambia de \(3/11\) a \(3/6\). Aquella fracción no era la traza de un
proyector estable. La corrección no elimina el sector: desplaza la lectura
al objeto homogéneo \(P_G\) de \(t=7\), donde rango y dimensión
pertenecen al mismo espacio.

\section{Consecuencia para una acción positiva}

La ley de acción puede usar ahora una única clase isométrica de proyectores
positivos en tres espacios de soporte:
\[
 P_G
 \xleftrightarrow{\ W_{GK}\ }
 P_3
 \xleftrightarrow{\ U_W\ }
 Q_3.
\]
La estrella \(R_{B_K}\) queda reservada para paridad y orientación. El
proyector \(P_G\) registra supervivencia; \(P_3\) registra la polarización
dodecafásica; \(Q_3\) registra su realización de incidencia en Witt. Esta
separación de funciones es la premisa matemática de la medida y de la ley
positiva de la sección siguiente. En el calibre fijo la cadena expresa
equivalencia isométrica de proyectores sin representación común. El
\cref{thm:recalibracion-hstar-witt} transporta la calibración completa y
convierte la cadena con \(\widetilde U_W\) y \(\widetilde Q_3\) en una
representación común de \(H_*\), sin identificar la estrella firmada con el
proyector positivo.
